% Part: incompleteness
% Chapter: introduction
% Section: overview

\documentclass[../../../include/open-logic-section]{subfiles}

\begin{document}

\olfileid[tr]{inc}{int}{ovr}

\olsection{Eksiklik Sonuçlarına Genel Bakış}

Hilbert, matematiğin !!{axiomatizable} bir kuramda biçimselleştirilebileceğini
ve bu kuramın !!{complete} ve !!{decidable} olduğunun ispatlanabileceğini
umuyordu. Üstelik sezgiciliğin meydan okumalarına karşı klasik matematiği
savunacak biçimde, bu kuramın tutarlılığını çok zayıf ``sonlu'' araçlarla
ispatlamayı amaçlıyordu. G\"odel'in eksiklik teoremleri bu hedeflere
ulaşılamayacağını gösterdi.

G\"odel'in birinci eksiklik teoremi, Russell ile Whitehead'in
\emph{Principia Mathematica}'sının bir sürümünün !!{complete} olmadığını
gösterdi. Ancak ispat gerçekte çok geneldi ve çok çeşitli kuramlara uygulanır.
Bu, matematiği bütünüyle yakalamayı yalnızca \emph{Principia Mathematica}'nın
başaramadığı değil, kabul edilebilir \emph{hiçbir} kuramın başaramayacağı
anlamına gelir. Eksiklik teoremlerinin uygulanması için yeterli olan kuram
özelliklerini ayırmak ve G\"odel'in ispatını yalnızca bu özelliklere bağlı
olacak biçimde genelleştirmek zaman aldı. Fakat artık birinci eksiklik
teoreminin, aritmetiğin $\Lang L_A$ dilindeki kuramlar için çok genel bir
sürümünü ifade edebiliriz.

\begin{thm}
$\Gamma$, $\Lang L_A$ içinde tutarlı ve !!{axiomatizable} bir kuram olup
bütün hesaplanabilir fonksiyonları !!{represents} ediyor ve bütün
!!{decidable} bağıntıları da temsil ediyorsa $\Gamma$ !!{complete} değildir.
\end{thm}

$\Gamma$'nın !!{complete} olmadığını söylemek, en az bir !!{sentence}~$!A$
için $\Gamma \Proves/ !A$ ve $\Gamma \Proves/ \lnot !A$ olduğunu söylemektir.
Böyle !!a{sentence}{}ye ($\Gamma)$ bakımından \emph{bağımsız} denir. Gerçekte
bağımsız tümceler bulunması gerektiğini oldukça çabuk ispatlayabiliriz.
Ancak G\"odel'in teorem ispatının gücü, böyle bağımsız bir !!{sentence}{}nin
\emph{belirli bir örneğini} ortaya koymasındadır. Bu ilgi çekici kuruluş,
$\Gamma$ için \emph{G\"odel tümcesi} denen !!a{sentence} üretir; onu
$!G_\Gamma$ ile gösteririz. Bu tümce ispatlanamaz, çünkü $\Gamma$ içinde $!G_\Gamma$, bizzat
$!G_\Gamma$'nın $\Gamma$ içinde ispatlanamaz olduğu iddiasına denktir.
Bunu \emph{yapıcı biçimde} yapar: $\Gamma$'nın bir aksiyomlaştırılması ve
!!{derivation} sisteminin bir betimi verildiğinde ispat, $!G_\Gamma$'yı
gerçekten yazmak için bir yöntem verir.

G\"odel'in ispatındaki kuruluş, $\Lang L_A$'nın terim ve !!{formula}s{}inin
özelliklerini ve bunlar üzerindeki işlemleri yine $\Lang L_A$ içinde dile
getirmenin bir yolunu bulmamızı gerektirir. Bunlara ``$!A$ !!a{sentence}{}dir'',
``$\delta$, $!A$'nın !!a{derivation}{}idir'' gibi özellikler ve
$\Subst{!A}{t}{x}$ gibi işlemler girer. Bu yol (a) simgeleri ve bunların
dizilerini (terimler, !!{formula}s, !!{derivation}s vb. bunlardır)
doğal sayılar (çünkü $\Lang L_A$ bunlardan söz edebilir) olarak
``kodlayarak'' söz konusu özellik ve bağıntıları dile getirmelidir. (b)
Bunu $\Gamma$'nın ilgili olguları ispatlayacağı biçimde yapmalıdır;
dolayısıyla bu özelliklerin doğal sayıların !!{decidable} özellikleriyle
kodlandığını ve işlemlerin doğal sayılar üzerindeki hesaplanabilir
fonksiyonlara karşılık geldiğini göstermeliyiz. Buna ``sözdizimin
aritmetikleştirilmesi'' denir.

Ancak sözdizimin nasıl aritmetikleştirilebileceğini araştırmadan önce,
$\Gamma$'nın ``yeterince güçlü'', yani bütün hesaplanabilir fonksiyonları
!!{represents} etmesi ve bütün !!{decidable} bağıntıları temsil etmesi
koşulunu ele alacağız.
Bu, ``hesaplanabilir''in kesin bir tanımını vermemizi gerektirir. Bu,
örneğin Turing makineleri modeli aracılığıyla veya genel amaçlı bir
programlama dilindeki programlarla hesaplanabilen fonksiyonlar olarak,
birkaç yolla yapılabilir. Bununla birlikte amacımız bu fonksiyon ve
bağıntıları $\Lang L_A$ dilindeki bir kuram içinde !!{represents} etmek
olduğundan, yalnızca sayısal fonksiyonların hesaplanabilirliğine ilişkin
yalın bir tanım seçmek en iyisidir. Bu, \emph{özyinelemeli fonksiyon}
kavramıdır. Bu yüzden önce özyinelemeli fonksiyonları tartışacağız. Ardından
$\Th{Q}$'nun bütün özyinelemeli fonksiyon ve bağıntıları zaten
!!{represents} ettiğini göstereceğiz. Böylece $\Th{Q}$ ve $\Th{PA}$ gibi
belirli kuramlara eksiklik teoremini uygulayabileceğiz; çünkü bunların
``yeterince güçlü'' kuram örnekleri olduğunu kurmuş olacağız.

Sözdizimin aritmetikleştirilmesinin son ürünü, !!a{formula}
$\OProv[\Gamma](x)$'tir; bu, !!{formula}s{}in sayılar olarak kodlanması
aracılığıyla $\Gamma$'nın aksiyomlarından ispatlanabilirliği dile getirir. Daha özel olarak $!A$,
$n$ sayısıyla kodlanmışsa ve $\Gamma \Proves !A$ ise
$\Gamma \Proves \Prov[\Gamma](\num{n})$ olur. $\Gamma$ için bu
``ispatlanabilirlik yüklemi'', $\Gamma$'nın tutarlılığını da belli bir
anlamda $\Lang L_A$ içinde !!a{sentence} olarak dile getirmemizi sağlar:
$\Gamma$ için ``tutarlılık ifadesi'', $n$'yi bir çelişkinin, örneğin
$\lfalse$'un kodu aldığımız $\lnot \Prov[\Gamma](\num{n})$ !!{sentence}{}si
olsun. İkinci eksiklik teoremi, tutarlı !!{axiomatizable} kuramların kendi
tutarlılık ifadelerini de ispatlamadığını söyler. Bu teoremin uygulanması
için gereken koşullar, kuramın bütün hesaplanabilir fonksiyonları ve
!!{decidable} bağıntıları temsil etmesinden biraz daha katıdır; fakat
$\Th{PA}$'nın bunları sağladığını göstereceğiz.

\end{document}
